\documentclass{jsarticle}
\usepackage[dvipdfmx,final]{graphicx}
\usepackage{amsmath}
\usepackage{amssymb}
\begin{document}
\title{誤差関数と3次元多様体、異次元への扉}
\author{}
\date{}
\maketitle

    自分がいると、方位が決まる。そうしたら、空間ができる。北と南が出来る。
    これは、磁気単極子が相対空間での根本的なアイデアらしい。
    無になると関係性がなくなり、磁気単極子もある上に、全て無でもある。
    無になると十二支縁起になる。宇宙と異次元の外が無でもある。
    運動量が決まると位置が決まる。位置だけだと宇宙が不確定性原理
    であり、異次元が合わさると運動量も位置も決まる。
    無になると位置も運動量も同時にわかる。
    このアイデアが3次元多様体のエントロピー値である。
    
    $$ \int^{\ll n} ||ds^2|| = \int\int\int 8 \pi G \left({p \over c^3} + 
    {V \over S}\right)dxdydz $$
    $$ = {\partial \over \partial f_n}{||ds^2||} $$
    \newcommand{\variint}{%
	\begin{picture}(15,30)
		\put(0,0){
			$ \iint $}
		\put(0,5){\line(15,0){15}}
	\end{picture} %
}

\newcommand{\variiint}{%
	\begin{picture}(20,30)
		\put(0,0){
			$ \iiint $}
		\put(0,5){\line(15,0){20}}
	\end{picture} %
}



 $$ {\partial \over \partial f}F(x) = \int \int \text{cohom}D_k(x)[I_m]  $$
 Cohomology element is constructed with isotopy of D-brane,
 and this isotopy of symmetry structure is sheaf of manifold,
 more over this brane of element is double integrate of projection.
 $$ {\partial \over \partial f}F = {{}^t \scalebox{1}[2]{\variint}}
 \text{cohom}D_k(x)^{\ll p} $$
 And global partial defferential equation is integrate of cut in cohomolgy,
 and this step of defferential D-brane of sheap value is varint of equals,
 inverse of horizen of cute of section value is reverse of integrate of 
 operators. Three dimension of varint of manifolds in operator of sheaps,
 and this step of times of value is equal with D-brane of equations.
 Global differential equation and integrate equation is operator of 
 global scales of horizen cut and add of cut equation are routs.

 $$ \oint r dx_m = \mathcal{O}(x,y) $$

 $$ {{}^t \scalebox{1}[2]{\variiint}}_{{D(\chi,x)}} 
 \text{Hom}[D^2 \psi]^{\ll p} \cong 
 \text{vol}\left({V \over S}\right) $$

 $$ {\partial \over \partial f}F(x) = {{}^t \scalebox{1}[2]{\variint}} 
 \text{cohom}D_k(x)^{\ll p} $$

 $$ {d \over df}F = (F)^{f^{'}}, \int F dx_m = (F)^{f}  $$
 $$ {d \over df}\int \int{1 \over (x \log x)^2}dx_m =
 \left(\int \int{1 \over (x \log x)^2}dx_m \right)^{f^{'}}  $$
 $$ f^{'} = (2x (\log x + \log (x + 1))) $$
  $$  \int \int F dx_m = \left({1 \over (x \log x)^2}\right)^{
	  (\int 2x (\log x + \log (x + 1))dx)} $$ 
   $$= e^{-f} $$
  $$ {d \over df}F = \left({1 \over (x \log x)^2}\right)^{
	  (2x (\log x + \log (x + 1))) } $$  
  $$= e^{f} $$
  $$ \log (x \log x) \ge 2(y \log y)^{1 \over 2} $$

  $$ \log (x \log x) \ge 2(\sqrt{y \log y}) $$

 
    $$ ||ds^2|| = 8 \pi G \left({p \over c^3} + {V \over S}\right) $$
    $$ \int ||ds^2||dx_m  = \int 8 \pi G \left({p \over c^3} 
    + {V \over S}\right)\mathrm{dvol}  $$
    $$ \int ||ds^2||dx_m  = \int 8 \pi G \left({p \over c^3} 
    + {V \over S}\right)dy_m $$
     $$ \int ||ds^2||dx_m  = \int {1 \over {8 \pi G \left({p \over c^3} 
    + {V \over S}\right)}}dx_m $$
    $$ {d \over df}F = m(x), \bigoplus \left({i \hbar}^{\nabla}\right)^{
	    \oplus L} $$
    $$ = \nabla_{i} \nabla_{j} \int \nabla f(x)d \eta $$

    $$ dx_m = {y \over {\log x}}, dy_m = {x \over \log x} $$
    $$ e^{-f}dV = dy_m = \mathrm{dvol} $$
    偏微分は加群分解と同じ計算式に行き着く。

   宇宙と異次元の誤差関数のエネルギー、
    $AdS_5$多様体がベータ関数となる値の列が、異次元への扉となっている。
    $$ \beta(p,q) = \text{誤差関数} + \text{Abel多様体} $$
    $$ = \text{$AdS_5$多様体} $$
    $$ = {d \over df}F + \int C dx_m = \int \Gamma(\gamma)^{'}dx_m $$
    ここで、アーベル多様体はEuler productである。
    ベータ関数の数列がわかると、ゼータ関数は無であるというのが、
    どういうことかが、物、物体に影ができて、ものが瞑想と同じであり、
    これから、風景がベータ関数の数値列に見えるらしい。
    映画のマトリックスの始めのシリーズでのネオの最後の覚醒と、
    この風景が電子の流れが、現実では、ものに影らしい。
    ものの生成物のフダの意味までもついているらしい。
    苫米地博士の私の秘密教えますの本と同じ能力らしい。

    この大域的微分多様体のガンマ関数が、複素力学系のマンデルブロ集合の
    プリズムと同じ構造の見方らしい。

    $$ ||ds^2|| = e^{-2\pi G ||\psi||}[\eta + \bar{h}(x)]dx^{\mu}dx^{\nu}
    + T^2 d^2 \psi $$
    $$ \beta(p,q) = \text{誤差関数} + \text{Abel多様体} $$
    $$ \int \mathrm{dvol} = \square \psi $$
    $$ \int \nabla \psi^2 d \nabla \psi = \square \psi $$
    $$ \text{expanding of universe} = \text{exist of value} $$
    $$ = \log (x \log x) = \square \psi $$
    $$ \text{freeze out of universe} = \text{reality of value} $$
    $$ = (y \log y)^{1 \over 2} = \nabla \psi $$

    All of value is constance of entropy, universe is freeze out
    constant, and other dimension is expanding into fifth dimension of
    inner.


      $$ x^n + y^n = z^n, \beta(p,q) = x^n + y^n - \delta(x) = z^n - \delta
      (x) $$
      $$ {d \over dt}g_{ij} = - 2 R_{ij}, {\cos x \over (\cos x)^{'}}
      \cdot (\sin x)^{'} = z_n, z^n = - 2 e^{x \log x} $$

      $$ z = e^{-f} + e^{f} - y $$
      $$ \beta(p,q) = e^{-f} + e^{f} $$
      相対性は、暗号解読と同じ仕組みの数式を表している。
      ここで言うと、yが暗号値である。テェックディジットと同じ仕組みを有
      している。

      $$ \square x = \int {f(x) \over {\nabla(R^{+}\cap E^{+})}}d \square x$$
      $$= \int {{\Delta f(x) \circ E^{+}} \over {\nabla(R^{+}\cap E^{+})}}
      \square x $$
      $$ \square x = \int {d(R \nabla E^{+}) \over {\nabla(R^{+} \cap E^{+})}}
      d \square x $$
      $$ = \int {{\nabla_{i} \nabla_{j} (R + E^{+}} \over {\nabla(R^{+}
      \cap  E^{+})}}d \square x $$
      $$ x^n + y^n = z^n $$
      $$ \mathrm{exp}(\nabla(R^{+}\cap E^{+}),\Delta(C \supset R)) $$
      $$ = \pi (R_1 \subset \nabla E^{+}) = \mathrm{rot}(E_1, 
      \mathrm{div}E_2) $$
      $$ xf(x) = F(x), s\Gamma(s) = \Gamma(s+1) $$
      $$ Q \nabla C^{+} = {d \over df}F(x) \nabla \int \delta(s)f(x)dx $$
      $$ E^{+}\nabla f = {{e^{x \log x} \nabla n! f(x)} \over {E(x)}} $$
      $$ {d \over df}F ~ F^{{f}^{'}} = e^{x \log x} $$
      $$ (C^{\nabla})^{\oplus Q} = e^{x \log x} $$
      $$ {d \over dt}g_{ij} = - 2 R_{ij}, {\cos x \over (\cos x)^{'}}
      \cdot (\sin x)^{'} = z_n, z^n = - 2 e^{x \log x} $$

      $$ R \nabla E^{+} = f(x) \nabla e^{x \log x}, {d \over df}F = F^{
	      {f}^{'}} = e^{x \log x}  $$
      南半球と単体（実数）の共通集合の偏微分した変数をどのようなF(x)かを
      $$ \int \delta(x)f(x)dx $$
      と同じく、単体積分した積分、共通集合の偏微分をどのくらいの微分変数
      を
      $$ \int \nabla \psi^2 d \nabla \psi = \square \psi $$
      と同じ、
      $$ \int dx = x + \text{C(Cは積分定数) }$$
      と原理は同じである。


\end{document}
